La proposición (6.1.1) admite el recíproco parcial siguiente:

\begin{proposition}[6.1.5]
\EGAReaderTarget{ega4.c04.prop.6-1-5}{IV.2.s6.1.5-fr}
Sean $A$, $B$ dos anillos locales noetherianos, $k$ el cuerpo residual de
$A$, $\varphi:A\to B$ un homomorfismo local y $M$ un $B$-módulo de tipo finito.
Se supone que:
\begin{enumerate}
\item[\textup{1\textsuperscript{o}}] $A$ es un anillo regular.
\item[\textup{2\textsuperscript{o}}] $M$ es un $B$-módulo de Cohen--Macaulay.
\item[\textup{3\textsuperscript{o}}] Se tiene
$\dim_B(M)=\dim(A)+\dim_{B\otimes_Ak}(M\otimes_Ak)$.
\end{enumerate}
Entonces $M$ es un $A$-módulo plano.
\end{proposition}

Procedamos por inducción sobre $n=\dim(A)$. Si $n=0$, $A$ es un cuerpo
puesto que es regular (\textbf 0, 17.1.4), y la afirmación es trivial. Supongamos
$n>0$; sea $\mathfrak m$ el ideal maximal de $A$, y sea
$x\in\mathfrak m$ un elemento que no pertenece a $\mathfrak m^2$; se sabe
entonces (\textbf 0, 17.1.8) que $A'=A/xA$ es regular y
$\dim(A')=\dim(A)\setmin 1$. Pongamos
\[
B'=B/xB,\qquad M'=M/xM=M\otimes_AA';
\]
se tiene, por tanto,
\[
B'\otimes_{A'}k=B\otimes_Ak,\qquad
M'\otimes_{A'}k=M\otimes_Ak
\]
y, en virtud de (5.5.1.2), se tiene
\[
\dim_{B'}(M')\leq\dim(A')+
\dim_{B'\otimes_{A'}k}(M'\otimes_Ak);
\]

\oldpage[IV-2]{137}

se concluye, pues, de (\textbf 0, 16.3.4) que se tiene
\[
\dim_B(M)\leq\dim_{B'}(M')+1
\leq\bigl(\dim(A')+
\dim_{B\otimes_Ak}(M\otimes_Ak)\bigr)+1
=\dim(A)+\dim_{B\otimes_Ak}(M\otimes_Ak).
\]

Como, por hipótesis, los miembros extremos de estas desigualdades son iguales,
se tiene necesariamente: (i) $\dim_{B'}(M')=\dim_B(M)\setmin 1$ y, como $M$ es por
hipótesis un $B$-módulo de Cohen--Macaulay, esto significa que $x$ es
$M$-\emph{regular} (\textbf 0, 16.1.9 y 16.5.5); (ii)
$\dim_{B'}(M')=\dim(A')+
\dim_{B'\otimes_{A'}k}(M'\otimes_{A'}k)$; como $M'$ es un $B'$-módulo de
Cohen--Macaulay en virtud de (i) y de (\textbf 0, 16.1.9 y 16.5.5), y $A'$ es regular,
la hipótesis de inducción prueba que $M'$ es un $A'$-módulo plano; se deduce
entonces de $(\mathbf{0}_{\mathrm{III}}, 10.2.7)$ que $M$ es un $A$-módulo plano,
puesto que $x$ es $M$-regular por (i).
